\chapter{Reference: spinners}
\label{ch:refman-spinner}

\begin{aside}
A spinner is one character cycled at a fixed rate. The whole
widget is a frame counter and a string of glyphs.
\end{aside}

\section{Working example}

\begin{lstlisting}[language=C++]
#include <maya/maya.hpp>
#include <chrono>
#include <string_view>

using namespace maya;
using namespace maya::dsl;

namespace {
constexpr std::string_view frames[] = {
    "|", "/", "-", "\\\\",
};
}

struct Spin {
    struct Model { int frame = 0; bool busy = true; };
    struct Tick {};
    struct Toggle {};
    struct Quit {};
    using Msg = std::variant<Tick, Toggle, Quit>;

    static Model init() { return {}; }

    static auto update(Model m, Msg msg)
        -> std::pair<Model, Cmd<Msg>>
    {
        return std::visit(overload{
            [&](Tick) {
                m.frame = (m.frame + 1) % 4;
                return std::pair{m, Cmd<Msg>{}};
            },
            [&](Toggle) {
                m.busy = !m.busy;
                return std::pair{m, Cmd<Msg>{}};
            },
            [](Quit) {
                return std::pair{Model{}, Cmd<Msg>::quit()};
            },
        }, msg);
    }

    static Element view(const Model& m) {
        auto glyph = std::string(frames[m.frame]);
        return v(
            h(
                when(m.busy,
                    text(glyph) | Fg<100, 200, 255>,
                    t<"*">      | Fg<100, 220, 120>),
                text(" "),
                text(m.busy ? "working" : "idle") | Bold
            ),
            blank_,
            t<"space toggle, q quit"> | Dim
        ) | pad<1> | border_<Round>;
    }

    static auto subscribe(const Model& m) -> Sub<Msg> {
        auto keys = key_map<Msg>({
            {'q', Quit{}},
            {' ', Toggle{}},
        });
        if (m.busy) {
            return Sub<Msg>::batch(std::move(keys),
                Sub<Msg>::every(std::chrono::milliseconds(80),
                                Tick{}));
        }
        return keys;
    }
};

int main() { run<Spin>({.title = "spinner"}); }
\end{lstlisting}

\section{Frame sets}

\begin{itemize}[leftmargin=*]
  \item ASCII: \code{|}, \code{/}, \code{-}, \code{\textbackslash}.
  \item Braille: \code{U+2800}--\code{U+28FF}, eight-frame
    cycles look smooth.
  \item Dots: \code{.}, \code{..}, \code{...} for a slower
    pulse.
\end{itemize}

\section{Recap}

\begin{recap}
\begin{itemize}[leftmargin=*]
  \item Frame counter modulo frame count.
  \item Tick at 60--120 ms for human-pleasant motion.
  \item Stop ticking when work finishes.
\end{itemize}
\end{recap}

\section*{Workshop}

\begin{enumerate}[leftmargin=*]
  \item Swap the ASCII frames for braille dots.
  \item Render three spinners with different periods.
  \item Show a green tick instead of the spinner once work
    is done.
\end{enumerate}
